% !TEX TS-program = pdflatex
% !TEX encoding = UTF-8 Unicode

% Example of the Memoir class, an alternative to the default LaTeX classes such as article and book, with many added features built into the class itself.

%\documentclass[12pt,a4paper]{memoir} % for a long document
\documentclass[12pt,a4paper,article]{memoir} % for a short document

\usepackage[utf8]{inputenc} % set input encoding to utf8

% Don't forget to read the Memoir manual: memman.pdf

%%% Examples of Memoir customization
%%% enable, disable or adjust these as desired


% This is from memman.pdf

%%% \maketitle CUSTOMISATION
% For more than trivial changes, you may as well do it yourself in a titlepage environment
\pretitle{\begin{center}\sffamily\huge\MakeUppercase}
\posttitle{\par\end{center}\vskip 0.5em}

%%% ToC (table of contents) APPEARANCE
\maxtocdepth{subsection} % include subsections
\renewcommand{\cftchapterpagefont}{}
\renewcommand{\cftchapterfont}{}     % no bold!

%%% HEADERS & FOOTERS
\pagestyle{ruled} % try also: empty , plain , headings , ruled , Ruled , companion

%%% CHAPTERS
\chapterstyle{hangnum} % try also: default , section , hangnum , companion , article, demo

\renewcommand{\chaptitlefont}{\Huge\sffamily\raggedright} % set sans serif chapter title font
\renewcommand{\chapnumfont}{\Huge\sffamily\raggedright} % set sans serif chapter number font

%%% SECTIONS
\hangsecnum % hang the section numbers into the margin to match \chapterstyle{hangnum}
\maxsecnumdepth{subsection} % number subsections

\setsecheadstyle{\Large\sffamily\raggedright} % set sans serif section font
\setsubsecheadstyle{\large\sffamily\raggedright} % set sans serif subsection font

%% END Memoir customization

\usepackage{tabularx}
\usepackage{graphicx}
\usepackage{amsmath}
\usepackage{amssymb}
\usepackage{amsthm}

\newtheorem{theoreme}{Théorème}[section]
\newtheorem{proposition}{Proposition}[section]
\newtheorem{definition}{Définition}[section]
\newtheorem{remarque}{Remarque}[section]


\title{Rapport de travail, Méthodes de bases réduites MS-13b}
\author{Arthur Timbeau-Périvier, Master AMS IP-Paris et Paris-Saclay}
\date{31/03/2026} % Delete this line to display the current date

%%% BEGIN DOCUMENT
\begin{document}
\maketitle

\newpage
\tableofcontents

\newpage
\section{Introduction}
Le fil conducteur du cours de méthode de base réduite est l'EDP suivante : 
\begin{equation}
    \label{fil_conduc}
    \begin{cases}
    -\text{div}(A(\mu)\nabla u(x)) = f(x,\mu) & \text{sur } \Omega \\
    u = 0 & \text{sur } \partial \Omega
    \end{cases}
\end{equation}


où $\mu \in \mathbb{R}^n$ représente un ensemble de paramètres propres au problème au quel on s'intéresse. Si on désire la résoudre numériquement, la voie naïve 
consiserait à la résoudre avec un schéma classique (comme une méthode d'éléments finis) pour un $\mu$ fixé, puis recommencer pour autant de jeu de paramètre pour lesquels on souhaite une solution. 
Cette idée est trop simpliste. Elle coûterait chère en temps de calcul, puisqu'il faudrait recalculer le schéma numérique pour chaque paramètre, 
mais surtout n'exploiterai pas un potentiel énorme gains de calcul intrasèque au problème posé. 

\vspace{0.3cm}

L'intuition est la suivante : connaître les solutions pour un certains nombre de jeu de paramètres suffit 
pour approcher la solution pour un jeu de paramètre encore non résolu. C'est l'idée derrière la principale méthode étudiée durant ce cours, 
la \textit{Proper Orthogonal Decomposition}, qui va consister en la réductions des dimensions du problème étudié. Une fois cette méthode mise en place, il va 
être question de la tester contre une solution haute fidélité pour voir à quel point on peut s'appuyer sur cette méthode, et étudier un cas particulier dans lequel 
elle ne fonctionne pas. 

\vspace{0.3cm}

Dans le TP1, on décrira plus en détail la POD, et on décrira une méthodes de résolution haute fidélité par rapport au quel on 
comparera la méthode de base réduite, à savoir la \textit{méthode des volumes finis}. Dans le TP2, on étudiera un cas dans lequel 
la méthode de base réduite est inéfficace et on essaiera de donner un sens à ce disfonctionnement. Et enfin, dans le TP3 et le TP4 on testera 
l'efficacité de la méthode de base réduite contre une autre méthode haute fidélité : la \textit{méthode des éléments finis} ainsi que l'étude de 
l'erreur a posteriori de cette dernière. 

\newpage

\section{TP1}
\subsection{La méthode des volumes finis}
La méthode des volumes finis repose sur la \textit{formule de Stokes}.

\begin{theoreme}
    Soit $u \in H^1(\Omega)$ où $\Omega$ est un ouvert de $\mathbb{R}^n$ à frontière lispschitzienne. Alors : 
    \[
        \int_{\Omega}\nabla u \, dx = \int_{\partial\Omega} u \cdot n\, d\sigma(x)
    \]
    où $\sigma$ défini une mesure sur $\partial \Omega$ et $n$ est une normale sortante de $\Omega$. 
\end{theoreme}

Ce résultat traduit une certaine relation entre une fonction à gradient intégrable et ses valeurs sur le bord du domaine (sous réserve que ce dernier soit suffisament régulier).

\vspace{0.3cm}

En intégrant \eqref{fil_conduc} sur $\Omega$, il vient :

\[
    \int_{\Omega}-\text{div}(A(\mu)\nabla u(x)) \, dx = \int_{\Omega}f(x,\mu) \, dx
\]

Maintenant, on applique la formule de Stokes au membre de gauche :
\[
    \int_{\Omega}-\text{div}(A(\mu)\nabla u(x)) \, dx = -\int_{\partial\Omega}A(\mu)\nabla u \cdot n\ \, d\sigma(x) 
\]

Ce qui permet d'écrire : 
\begin{equation}
    \label{main_rel_VF}
    -\int_{\partial\Omega}A(\mu)\nabla u \cdot n\ \, d\sigma(x)  = \int_{\Omega}f(x,\mu) \, dx
\end{equation}

\vspace{0.3cm}

On considère $\mathcal{M}_h$ un maillage du domaine $\Omega$ où $h$ est la longueur de la plus grande arête du maillage. On dénote 
par $\mathcal{F}$ l'ensemble des arêtes de $\mathcal{M}_h$. Pour chaque élément $K$ de $\mathcal{M}_h$, on réécrit la relation \eqref{main_rel_VF} :

\[
    -\int_{K}A(\mu)\nabla u \cdot n \, d\sigma(x)  = \int_{K}f(x,\mu) \, dx
\]

Ce qui peut encore se redécomposer en somme d'intégrale sur les arêtes $e$ de $K$ :

\[
    \sum_{e \in \mathcal{F_K}}-\int_{e}A(\mu)\nabla u \cdot n_{e,K} \, d\sigma(x)  = \int_{K}f(x,\mu) \, dx
\]
où $n_{e,K}$ désigne l'arête sortante de l'arête $e$.

L'objectif est de construire, à partir de cette dernière formulation, un schéma numérique qui va nous permettre de calculer 
une solution approchée de $u$. Premièrement, on note 

\begin{align*}
 &\frac{1}{|K|}\int_{K}f(x,\mu) \, dx = f_K \\
 &-\int_{e}A(\mu)\nabla u \cdot n_{e, K} \, d\sigma(x)
\end{align*}

Ensuite, on considère $L$ un élément de $\mathcal{M}_h$ qui partage une arête $e$ avec $K$. On a 
\[
    \int_{e}A(\mu)\nabla u \cdot n_{e,L} \, d\sigma(x) = -  \int_{e}A(\mu)\nabla u \cdot n_{e,K} \, d\sigma(x) 
\]

et on fait l'approximation 

\[
    A(\mu)\nabla u \approx A^{K}(\mu)\frac{(u_e - u_K)}{d(x_K, e)}
\]
où $x_K$ est le centre de l'élément $K$ et $u_K = \frac{1}{|K|}\int_{K}u \, dx$. On a donc

\begin{align*}
    &F_{K,e} \approx A^{K}(\mu)\frac{(u_e - u_K)}{d(x_K, e)} \sigma(e) \\
    &F_{L,e} = -F_{K_e} 
\end{align*}

Ce qui, après quelques calculs, aboutit à

\[
F_{K,e} = \frac{A^{K}(\mu)A^{L}(\mu)\sigma(e)}{A^{K}(\mu)d(x_K, e) + A^{L}(\mu)d(x_L, e)}(u_K - u_L)
\]
\vspace{0.3cm}

Ce qui nous donne le schéma que l'on souhaitait, car avec ce schéma on peut maintenant construire la matrice en parcourant chaque élément du maillage et
en réajoutant de proche en proche les contributions de chaque arête d'un élément. Le second membre de construit à l'aide d'une méthode de quadrature
pour approcher la valeur de $f_K$. 

\vspace{0.3cm}

C'est d'ailleurs ainsi que l'algorithme construit la matrice dans
 le code. On calcule les contributions pour les arêtes horizontales, 
on ajoute également les contributions des arêtes horizontales et on traite de manière spéciale les flux sur le bord du maillage. Il s'agit plus
ou moins de la même formule que ci-dessus. Cela dépend de où ce trouve le noeud : s'il n'est pas sur le bord, alors on prend
\[
    F_{K,e} = -A^{K}(\mu)\frac{(u_e - u_K)}{d(x_K, e)}\sigma(e)
\]

Et ce qui aboutit, grâce à la condition de Dirichlet sur le bord, 

\[
    F_{K,e} = A^{K}(\mu)\frac{u_K}{d(x_K, e)}\sigma(e)
\]

Si le noeud se trouve sur l'arête, alors on a 
\[
    F_{K,e} = 0
\]

soit, pour $g \in L^2(\Omega)$ une condition de bord quelconque
\[
u_K = \frac{1}{\sigma(e)}\int_{e}g \, d\sigma(x)
\] 
Et dans notre cas : 
\[
    u_K = 0
\]

Une fois que la matrice des volumes finis et que le second membre sont assemblés, il n'y a plus qu'à résoudre le système.
On remarque que la méthode des volumes finis repose grandement sur la conservation du flux à l'interface
entre 2 éléments.

\vspace{0.3cm}

Très adapté au problèmes hyperboliques, la méthode des volumes finis constitue une méthode de résolution haute fidélité
au même titre que la méthode des éléments finis. Ceci étant dis, comme expliqué plus haut, on ne peut pas
se permettre d'utiliser cette méthode à chaque fois que le paramètre $\mu$ change. 

\vspace{0.3cm}

L'idéal serait de disposer d'une formule exacte permettant de passer de la solution HF pour un $\mu_0 \in \mathbb{R}^n $ à
n'importe quel $\mu$ : 
\[
 u(x, \mu) = f(u(x,\mu_0), \mu)
\]

Disposer de la formule exacte de $f$ serai inespéré, mais supposer que cette relation existe est possible ne l'est pas. 
D'un point de vue d'un statisticien, on peut voir cette relation comme une famille d'expériences aléatoires dépendante
d'un paramètre et dont on voudrait estimer la loi. On va se baser sur les observations passées pour identifier une approximation de $f$
puis inférer sur des nouveaux paramètres. C'est le point de départ de la \textit{POD}, connue chez les statisticiens sous le nom de 
\textit{Analyse en Composante Principale}.

\subsection{POD et interprétation}

\vspace{0.2cm}

Le plan est le suivant : tout d'abord, calculer pour plusieurs jeu de paramètre $\mu$ la solution HF avec la méthode des volumes finis (que l'on appelle \textit{Snapshot en $\mu$}). 

\vspace{0.2cm}

Ensuite, on calcule
la covariance entre chacun des snapshot. Cette matrice de covariance va être, par construction, une matrice symétrique, donc diagonalisable dans $\mathbb{R}$ dans une base unitaire et à valeurs propres réelles et positive.

\vspace{0.2cm}

On choisit les sous-espaces propres associés au plus grande valeurs propres et on récupère cette nouvelle famille libre et génératrice (ce n'est pas une base !) dans laquelle 
on effectuera tout nos calculs. 

\vspace{0.3cm}

Ce que l'on vient de décrire est la \textit{Décomposition en Valeur Singulière} de la matrice des snapshots. L'idée de réduction 
apparaît lorsque l'on choisit, parmis les valeurs propres de la matrice de corrélation, celle étant les plus importantes pour décrire la relation entre deux snapshots.

\vspace{0.3cm}

Dans les cas où cette méthode est viable, 2 gains sont réalisés à sa mise en place. Primo : lorsque que le maillage est fin, 
il est beaucoup plus rapide d'effectuer ces calculs plutôt que de recalculer une solution HF. Secondo : on peut calculer \textit{offline} la famille orthonormale réduite
pour n'avoir qu'à inverser un système pour approcher la solution de \eqref{fil_conduc} (calcul \textit{online}). L'étude de performance
de cette méthode sera fait dans la section suivante.  

\vspace{0.3cm}

Pour s'assurer que cette idée tient la route, mettons un pied dans les détails. Un vecteur $\phi$ de la base réduite 
minimise l'erreur quadratique entre lui même et tous les snapshots (ou maximise la variance entre lui et chacun des snapshots), ce qui se formule comme suit :
\begin{align*}
    \begin{cases}
        &\text{Trouver}\; \phi \; \text{tel que} \\
        &||\phi|| = 1 \; \text{et} \; \phi \; \text{minimise} \; \mathbb{E}_\mu \left[||u - \left(u, \phi\right)_{L^2} \phi ||^2 \right]
    \end{cases}
\end{align*}

Et on utilisera plutôt la formulation "duale"

\begin{equation}
    \label{opt_pb}
    \begin{cases}
    &\text{Trouver}\; \phi \; \text{tel que} \\
    &||\phi|| = 1 \; \text{et} \; \phi \; \text{maximise} \; \mathbb{E}_\mu \left[|\left(u, \phi\right)_{L^2}|^2 \right]
    \end{cases}
\end{equation}

On introduit la forme quadratrique 
\begin{align*}
    Q :\; &L^2(\Omega) \longrightarrow \mathbb{R}_+ \\
        &(\phi, \psi) \; \mapsto \mathbb{E}_\mu \left[|\left(u, \phi\right)_{L^2}||\left(u, \psi\right)_{L^2}| \right]
\end{align*}

Cette forme quadratique est fermée (bornée inférieurement par 0, symétrique et l'espace $(L^2, \langle \cdot , \cdot \rangle_Q)$ est un espace de Hilbert),  
Par conséquent, l'opérateur $C$ associé à $Q$ est auto-adjoint sur $L^2$ et est bornée inférieurement. Identifions l'opérateur $C$ :

\begin{align*}
    Q(\phi , \psi) &= \mathbb{E}_\mu \left[|\left(u, \phi\right)_{L^2}||\left(u, \psi\right)_{L^2}| \right] \\
                   &= \int\int(u(\mu), \phi)u(\mu) \psi \; dx \, d\mu \\
                   &= \int \int(u(\mu), \phi)u(\mu) \, d\mu \; dx \\
                   &= \int \mathbb{E}_\mu \left[(u(\mu),\phi)u\right] \psi \, dx \\
                   &= \left(C\phi, \psi\right)
\end{align*}

où on a utilisé le théorème de Fubini. Ainsi, l'opérateur $C$ associé à $Q$ est définit par $C\phi = \mathbb{E}_\mu \left[(u(\mu),\phi)u\right]$. L'opérateur $C$ est compact : pour $(\phi_n)_{n\in\mathbb{N}}$ une famille orthonormale de 
$L^2$, 2 applications successives du théorème de convergence dominée puis la faible convergence vers 0 de la famille montre que 
\[
||C\phi_n||_{L^2} \underset{n \to \infty}{\longrightarrow} 0
\]

Mieux encore : $C$ est un opérateur à noyau intégral :
\[
C\phi(x) = \int K(x,y)\phi(y) \,dy
\]
Avec $K(x,y) = {E}_{\mu}\left[u(\mu, x)u(\mu, y)\right] \in L^2(\Omega \times \Omega)$. C'est donc un opérateur de Hilbert-Schmidt. De là, on peut tirer 2 conclusions :
\begin{itemize}
    \item[(I)] Par le théorème spectral pour les opérateurs compacts, il existe une base hilbertienne associée
                           à une suite $(\lambda_n)_{n\in\mathbb{N}}$ positive qui tend vers 0 dans laquelle $C$ est diagonal
    \item[(II)] Il existe une suite d'opérateur de rang fini qui converge dans l'espace dual vers $C$.
\end{itemize}

La première justfie l'existence de solution des problèmes spectraux associées à $C$, ce qui permet d'établir l'existence de solution à \eqref{opt_pb}.
La deuxième justifie le choix d'une suite d'opérateur de rang finis pour approximer numériquement l'opérateur $C$. Le fait qu'il soit de Hilbert-Schmidt
nous aiguille sur le choix de la suite d'opérateur. Maintenant que le problème en dimension infinie est bien posé, tâchons de le traduire en dimension finie. 
\newpage

Le noyau peut s'approcher par la moyenne empirique sur l'ensemble des snapshots déjà calculé : 
\[
K(x,y) \approx \frac{1}{N_T}\sum_{k = 1}^{N_T}u(x,\mu_k)u(y,\mu_k)
\]

Donc 

\[
C\phi \approx \frac{1}{N_T}\sum_{k = 1}^{N_T}u(x,\mu_k) \int_{\Omega}u(y,\mu_k)\phi(y)\, dy =  \frac{1}{N_T}\sum_{k = 1}^{N_T}u(x,\mu_k) (u(\mu),\phi)_{L^2(\Omega)}
\]

Et donc pour $(\lambda, \phi)$ une paire valeur/vecteur propre de $C$, on obtient
\[
\lambda \phi \approx  \frac{1}{N_T}\sum_{k = 1}^{N_T}u(x,\mu_k) (u(\mu),\phi)_{L^2(\Omega)}
\]

En décomposant $\phi$ sur la base des snapshots : $\phi(x) = \sum_{l = 1}^{N_T}\alpha_l u(x,\mu_l)$, le membre de droite devient
\[
\frac{1}{N_T}\sum_{k = 1}^{N_T}u(x,\mu_k) (u(\mu_k),\phi)_{L^2(\Omega)} = \frac{1}{N_T}\sum_{k = 1}^{N_T}u(x,\mu_k)\sum_{l = 1}^{N_T}\alpha_l(u(\mu_k), u(\mu_l))_{L^2(\Omega)}
\]

On y est presque, le dernier ingrédient à ajouter consiste à approcher le produit scalaire de $L^2(\Omega)$ par le produit scalaire euclidien
normalisé par le volume du domaine $\Omega$ (ce qui signifie discrétiser le domaine) : 
\[
    (u(\mu_k), u(\mu_l))_{L^2(\Omega)} \approx |\Omega|\sum_{i = 1}^{n}u(x_i, \mu_k)u(x_i, \mu_l) = |\Omega|\text{Cov}\left[(u(\mu_k)_i)_{i\leq n}, (u(\mu_l)_i)_{i\leq n}\right]
\]

En posant $\mathbb{C} = \left(\text{Cov}\left[(u(\mu_k)_i)_{i\leq n}, (u(\mu_l)_i)_{i\leq n}\right]\right)_{1 \leq k,l\leq N_T}$ et $\Phi = (\alpha_i)_{i\leq N_T}$, on reformule le problème
spectral sur $L^2(\Omega)$ en un problème spectral sur $\mathbb{R}^{N_T}$ :

\begin{equation}
    \mathbb{C}\Phi = \lambda\Phi
\end{equation} 
Et par construction, $\mathbb{C}$ est bien symétrique et la base réduite calculé est bien orthonormale, comme annoncé. La matrice $\mathbb{C}$ est appelé \textit{Matrice de corrélation}, terminologie issue 
de l'approche probabiliste du problème. 

\vspace{0.3cm}

Une fois que cette matrice est construite, on la diagonalise. On récupère les valeurs propres les plus importantes 
(par plus importante, comprendre les plus grandes), étant la discrétisée d'un opérateur compact, on est sûr que ces valeurs propres forment
une suite tendant vers 0 à mesure que le nombre de snapshots augmente. De là, on normalise les vecteurs propres par $\sqrt{\lambda}$ et on obtient à la fin la 
fameuse base réduite tant convoitée. 

\vspace{0.3cm}

Remarquons que la construction de $\mathbb{C}$ repose sur la quantité de snapshots précalculé. Ce qui signifie que la base réduite
peut être calculée en amont, puis réutilisé pour résoudre \eqref{fil_conduc} pour un nouveau jeu de paramètre $\mu$. 
La résolution haute fidélité mène à la résolution d'un système
\[
\mathbb{A}(\mu)U = F
\]
avec $\mathbb{A}(\mu) \in \mathcal{M}_{n\times n}(\mathbb{R})$, $U, F \in \mathbb{R}^n$ et $n$ le nombre de noeud sur le 
maillage. Pour résoudre le problème pour un nouveau $\mu$, on se place simplement dans la base réduite (orthonormale) $P$ :
\[
PP^T\mathbb{A}(\mu)PP^TU = F(\mu)
\]   

et alors pour $\Phi = P^TU$, on peut réécrire le système comme suit : 

\[
P^T\mathbb{A}(\mu)P\Phi = P^TF(\mu)
\] 

et la fonction $f$ introduit en amont décrivant la relation entre les snapshots (dont on rêve d'avoir une formule explicite) 
est approchée par la matrice $P^T\mathbb{A}(\mu)P$.  

\vspace{0.3cm}

En faisant une hypothèse supplémentaire sur la relation en $\mu$ et $\mathbb{A(\mu)}$ et $F(\mu)$, à savoir le fait
que la matrice dépend de manière affine de $\mu$, on peut encore accélérer la procédure.

\newpage


\subsection{Résultats numériques et erreur}
L'idée 

\newpage

\section{TP2}
\subsection{Epaisseur de Kolmogorov}
\subsection{Cas particulier}

\newpage

\section{TP3 et TP4}
\subsection{Base réduite contre les éléments finis}
\subsection{Erreur a posteriori}


\end{document}
